% =====================================================================
\chapter{AI Model Knowledge and Technical Understanding}
\label{sec:models}
% =====================================================================
Numbers in this chapter are from late November 2026. Some model rankings come from summary websites that often repeat
companies' own numbers; I treat those as lower-confidence and rely mainly on an independent benchmark where every model
runs in the same test setup \citep{valstb4}. Numbers from my own prototype are marked with a decision number.

\section{The main types of AI models, and where each fits in my project}

\begin{table}[H]
\centering
\small
\caption{Model types, examples, and their role in this project.}
\label{tab:categories}
\begin{tabularx}{\linewidth}{@{}P{2.6cm}P{3.2cm}P{3.3cm}L@{}}
\toprule
\textbf{Type} & \textbf{What it does} & \textbf{Examples (late 2026)} & \textbf{Where it fits in my project} \\
\midrule
Large language models (LLMs) & Read and write text, reason, follow instructions &
  Claude Opus/Sonnet 5.5, GPT-6.x, Gemini 4; open: Kimi K3, GLM-5.x, DeepSeek V4, Qwen3.x &
  The \textbf{AI second reader}. In the prototype this is Qwen3-VL-8B-Instruct used with text only (the same model that
  reads the pages), so only one model sits on the GPU. \\
Multimodal (vision-language) & Look at an image and answer in text & Qwen3-VL, MiniCPM-V, Gemini &
  \textbf{Reads the scanned scripts}: answers, ticks, red marks (main model: Qwen3-VL-8B). \\
Speech-to-text / text-to-speech & Speech to text, text to speech & Whisper-family open models; open TTS models &
  Not used yet. Later: explain a re-mark decision to parents \textbf{by voice, in Urdu}. \\
Computer vision & Find, cut out or classify things in images & YOLO, SAM, OpenCV &
  \textbf{Red-ink separation} (a colour rule). Next step: a detector that finds and crops every red mark. \\
Image and video generation & Create images or video from a prompt & Diffusion models &
  \textbf{Not for marking} (no content may be invented). Possible use: practice data (synthetic marked scripts), which
  must be checked against real scripts. \\
Embedding and reranking & Find similar texts and sort them by relevance & Open embedding and reranker models &
  Find \textbf{similar answers already marked by senior examiners} as examples for the second reader; spot copied answers. \\
Small language models & Small, cheap, fast models for normal hardware & Qwen3-VL-4B, GLM-OCR (0.9B) &
  Tested in my bake-off. The 4B model was \textbf{worse and also slower} than the 8B \dec{D016}. \\
AI coding and agent models & Write code, use tools, finish multi-step tasks & Claude Code, Codex, Gemini agents &
  I used Claude as a coding and research assistant to build the prototype (Section~\ref{sec:ai-help}). \\
\bottomrule
\end{tabularx}
\end{table}

% ---------------------------------------------------------------------
\section{Question 1: choosing a model for a tool-using, multi-step agent}

\paragraph{What an agent model has to be good at.} Being ``smart'' is not enough. The model has to plan a task over many
steps, call tools with the right inputs, read what the tool returns and recover when something goes wrong, remember what
it has done during a long task, and know when to stop or ask for help. So I judge models on agent tests (Terminal-Bench,
SWE-bench Pro, OSWorld, tool-use tests), on cost and time per \emph{finished} task, and on how they do on my own tasks.

\begin{table}[H]
\centering
\small
\caption{Shortlist for an agent model.}
\label{tab:shortlist}
\begin{tabularx}{\linewidth}{@{}P{2.2cm}P{3.6cm}LL@{}}
\toprule
\textbf{Family} & \textbf{Models} & \textbf{Strengths} & \textbf{Weaknesses} \\
\midrule
Claude & Opus 5.5 (planning), Sonnet 5.5 (main worker), Haiku 5.5 (helpers) &
  Top of the independent same-setup agent test, about 65\% on Terminal-Bench 4.0 \citep{valstb4}; strong on long tool use &
  Among the more expensive per task; closed, API only \\
GPT & GPT-6 Astra (hard reasoning), GPT-6.1 Sol (value) &
  Very strong reasoning and maths; Sol is cheap per task \citep{geotoolbox2026,aitoolsreview2026} &
  Astra is expensive; a little behind on the same agent test \citep{valstb4} \\
Gemini & Gemini 4 Carbon / Argon &
  Fastest per task in \citep{valstb4}; leads some app-agent rankings \citep{geotoolbox2026}; built for images, audio, video &
  Ranking changes a lot between leaderboards; closed \\
Open-weight & Kimi K3/K3.1 Pro, GLM-5.x, DeepSeek V4.x, Qwen3.x (27B fits one GPU) &
  Can run on our own machines; Kimi K3.1 Pro reached about 98\% of the top score at about 38\% of the cost
  \citep{valstb4,morphopen2026} &
  The largest need many-GPU servers; more work to operate; licences differ \\
\bottomrule
\end{tabularx}
\end{table}

\paragraph{How I would compare them.} First, \textbf{leaderboards disagree}: one ranks Claude first for agents, another
ranks Gemini 4 Argon first \citep{geotoolbox2026,benchlmagents2026}. Companies' own numbers can also look better than
reality: on SWE-bench Pro's private set of unseen code, the same model scored 59.1\% on the public set but 43.4\% on the
private set \citep{morphswepro}. So I would build a small test from my own work (20--50 real tasks with a checkable end
result), run every shortlisted model in the \emph{same} setup, and measure task success, tool-call errors, recovery after
a failed tool call, cost per successful task, time per task, and variation between runs. This is exactly what I did in
the prototype's bake-off \dec{D011,D016}.

\paragraph{From my own experience: choosing a local agent model.} This semester I faced the same question in my
course project on multi-agent LLM software teams, where three agents write code, run tests and message each other
through tools over many steps \citep{agenticlog}. There was no budget for paid APIs, so the agents had to run locally
on my own GPU (RTX 5060 Ti, 16\,GB). Before running anything, I fixed the selection rule: pick the main model by
\emph{tool-call failures}, \emph{whether the agents actually update the code after a change} (the change's decision test),
and \emph{minutes per run} \citep{agentic0043}. Then I ran the same smoke test on each candidate (Table~\ref{tab:q1local}).

\begin{table}[H]
\centering
\small
\caption[Smoke test of local agent models in my multi-agent course project (same project, same memory setting, one 5-episode
run each).]{Smoke test of local agent models in my multi-agent course project (same project, same memory setting, one 5-episode
run each) \citep{agentic0049,agenticlog}.}
\label{tab:q1local}
\begin{tabular}{@{}lccc@{}}
\toprule
 & \textbf{qwen3-coder:30b} & \textbf{gpt-oss:20b} & \textbf{gpt-oss:20b} \\
 & & (reasoning sent back) & (not sent back) \\
\midrule
Rejected tool calls (wrong tool name or arguments) & \textbf{0 / 91} & 65 / 105 & 99 / 118 \\
Decision tests passed after a change & \textbf{3 / 4} & 0 / 4 & 0 / 4 \\
Minutes per run & 11.4 & 10.0 & 7.3 \\
\bottomrule
\end{tabular}
\end{table}

\textbf{I chose qwen3-coder:30b} (a mixture-of-experts model with about 3B active parameters) because it was the only
candidate whose agents actually acted on changes \citep{agentic0049}. gpt-oss:20b was slightly faster, but in my setup its
tool calls broke down: it invented tool names and arguments, wrote calls as plain text, and repeated the same rejected call
up to 90 times, making only one commit in ten episodes. Earlier, the smaller qwen3:14b ran but did not build the specified
code interface and fell into repetition loops with greedy decoding \citep{agentic0038}. As an upper reference, a small paid
Claude model (Haiku 5.5) built the specified interface correctly in the same project \citep{agenticlog}. Three lessons from
this shaped my answer to this question:
\begin{itemize}
  \item \textbf{Benchmarks and model size are not enough.} A model that looks good on paper can fail completely on
        \emph{my} tool format; only a test on my own task showed it.
  \item \textbf{The tool layer matters as much as the model.} At first qwen3-coder also wrote every tool call as text,
        because the client library did not detect its tool support. After a small fix it made zero bad calls
        \citep{agentic0046}. Without checking this, I would have rejected the best model.
  \item \textbf{Follow the model card's settings.} Greedy decoding caused endless repetition in qwen3:14b, as its model card
        warns \citep{agentic0038}.
\end{itemize}

\paragraph{My choice.} \textbf{Claude Sonnet 5.5 as the main agent, Opus 5.5 only for the hardest planning steps, and a
cheap, fast model (Haiku 5.5 or a Gemini Flash model) for simple helper tasks.} The Claude models lead the independent
same-setup agent test, and Sonnet is within about one point of the top at lower cost and time \citep{valstb4}. Sending
easy work to small models usually saves more than switching companies. \textbf{If the data must stay inside the building},
as with exam scripts, I would use an open-weight model on our own GPUs (Kimi K3.x or GLM-5.x on a multi-GPU server, or
qwen3-coder:30b on a single 16\,GB GPU, the model I chose by testing in my own project above). This choice is temporary; I would re-run my own test every few months.
\emph{Note on bias: this answer was drafted with help from Claude, a Claude model; to reduce bias it relies on an
independent benchmark and shows where GPT and Gemini lead} \dec{D014}.

% ---------------------------------------------------------------------
\section{Question 2: two models give similar results, but one is faster and cheaper}

\paragraph{First, check that the results really are similar.} An average can hide big differences. Before I call two
models equal, I check: accuracy on my own data with an error margin (with 50 test items, a 2--3 point gap is usually
chance); \emph{where} each model fails, not just the average; consistency across repeated runs; whether the model knows
when it is unsure; speed (typical response time and the slowest 5\% of responses, plus behaviour under load); and
\textbf{cost per successful task}, including retries and human fixes, not cost per token.

\paragraph{Tests I would use.} Public tests only to shortlist (Terminal-Bench and SWE-bench Pro for agents, OCR and
handwriting tests for reading \citep{wildhandbench2026,omnihand2026}, MMMU for images). My own test set makes the final
decision. Example: on the same agent test, Kimi K3.1 Pro scored 63.64\% for \$4.57 and Claude Opus 5.5 scored 65.15\% for
\$12.04 \citep{valstb4}. The scores are almost the same, but the cost differs by about 2.6x.

\paragraph{Testing on real use.} (1) An offline test on real past cases, including hard and rare ones. (2) Shadow mode:
run the cheap model next to the current one on real requests without users seeing its answers. (3) A small rollout to
a share of real traffic with close monitoring and quick rollback. (4) Study the cases where the two models disagree.

\paragraph{When I would choose the slower, more expensive model.}
\[
\text{real cost per task} = \text{running cost} + P(\text{mistake}) \times \text{cost of that mistake}
\]
I choose the expensive model when mistakes are expensive (in my audit, a missed under-marking can cost a student a college
seat), for long agent tasks where small per-step errors add up, when the cheap model fails on rare but important cases,
or when volume is small so the saving is tiny. Otherwise the cheaper model wins. A good middle way is \textbf{routing}:
easy cases to the cheap model, hard or unsure cases to the big one.

\begin{table}[H]
\centering
\caption[A real example from my bake-off (5 scripts, same method A).]{A real example from my bake-off (5 scripts, same method A) \dec{D016}.}
\label{tab:q2example}
\begin{tabular}{@{}lcccc@{}}
\toprule
\textbf{Model} & \textbf{Questions found} & \textbf{Student letters} & \textbf{Red numbers} & \textbf{s/page} \\
\midrule
Qwen3-VL-8B-Instruct & 95.7\% & 100\% & 12.9\% & 8.7 \\
Qwen3-VL-4B-Instruct (smaller) & 58.3\% & 40\% & 11.3\% & 15.6 \\
\bottomrule
\end{tabular}
\end{table}
I expected the smaller model to be faster and only a little worse. On my pages it was \textbf{both slower and much worse},
probably because it wrote longer, messier answers for full pages. The lesson: measure speed and quality on your own task;
model size and price per token do not tell you the cost of a \emph{successful} task.

% ---------------------------------------------------------------------
\section{Question 3: what I would build with an in-house GPU server}

\paragraph{1. What and why.} A \textbf{``second reader'' for board exam marking at provincial or national scale}, the full
version of my prototype. It reads every scanned answer script (English and Urdu, handwritten), checks the examiner's marks
against themselves and against an independent AI reading, and gives each board a \textbf{list of scripts ranked by risk},
so the human re-checking boards already pay for goes where errors are most likely. It also gives students a fair,
evidence-based rechecking. The AI never gives the final mark. The problem is described in Chapter~\ref{sec:problem}.

\paragraph{2. Why an in-house GPU.} \emph{Secrecy:} a government board cannot send millions of children's scripts to a
foreign AI company. \emph{Volume and cost:} about 80 million pages a year (my estimate: about 20 million scripts $\times$
about 4 pages); API fees grow with every page, while owned hardware becomes cheaper per page as volume grows.
\emph{Training on our own data:} ready-made models do not handle Urdu handwriting and Pakistani answer formats well, and
board scripts cannot leave the building. \emph{Deadlines:} results come out on fixed dates and cannot depend on the
internet or API limits.

\paragraph{3. Models and technologies.} A vision-language page reader (Qwen3-VL-8B class won my bake-off: 99\% of student
answers but only 27\% of examiner red numbers read correctly \dec{D018}, so a larger or fine-tuned model is needed for red
marks), plus a small OCR model for printed parts; classic computer vision (red-ink separation, straightening pages,
splitting questions); an open-weight language model (around Qwen3.x 27B) as the second reader; embeddings and a reranker
to retrieve similar answers already marked by senior examiners; and, later, speech output for parents.

\paragraph{4. Existing, fine-tuned or from scratch?} \textbf{Fine-tune (LoRA) existing open models.} Training from scratch
needs far more data and compute and is not needed, because current models already read handwriting reasonably well.
But they still struggle with handwriting \citep{wildhandbench2026}, and none were trained on Urdu answer scripts with red
marks. LoRA (a cheap way to adjust a model) on a few thousand scripts labelled by senior examiners is affordable, stays
inside the building, and can be repeated every exam season.

\paragraph{5. Architecture.} Figure~\ref{fig:q3arch} shows the full system. Models are served with vLLM, which handles many
requests at once. Reviewer decisions flow back as new training data.

\begin{figure}[H]
\centering
\begin{tikzpicture}[node distance=4mm and 4mm, every node/.style={font=\footnotesize},
  box/.style={draw=accent,fill=soft,rounded corners=2pt,align=center,minimum height=9mm,text width=2.25cm},
  arr/.style={-{Stealth[length=2mm]},accent,thick}]
\node[box] (scan) {Scanner +\\secure storage};
\node[box,right=of scan] (prep) {Page prep:\\straighten, split\\red / student ink};
\node[box,right=of prep] (read) {Page reader\\(vision-language\\model)};
\node[box,right=of read] (rules) {Rule checks\\(totals, max,\\unmarked)};
\node[box,below=8mm of rules] (second) {AI second reader\\+ retrieved\\examples};
\node[box,left=of second] (risk) {Risk score\\and ranking};
\node[box,left=of risk] (queue) {Review screen for\\senior examiners};
\node[box,left=of queue] (train) {Their decisions\\become training\\data + audit log};
\draw[arr] (scan)--(prep); \draw[arr] (prep)--(read); \draw[arr] (read)--(rules);
\draw[arr] (rules)--(second); \draw[arr] (second)--(risk); \draw[arr] (risk)--(queue); \draw[arr] (queue)--(train);
\draw[arr,dashed] (train.north)--++(0,3mm)-|(read.south);
\end{tikzpicture}
\caption{Proposed in-house system. Everything runs inside the board's building; no internet is needed.}
\label{fig:q3arch}
\end{figure}

\paragraph{6. Hardware and infrastructure.} \emph{Running the models (an estimate from my prototype):} one consumer GPU
(RTX 5060 Ti, 16\,GB), one request at a time, read about \textbf{300 pages per hour} (median 48\,s per 4-page script
\dec{D021}). About 80 million pages $\div$ 300 $\approx$ 267{,}000 GPU-hours; in a 6-week marking period running day and night
(about 1{,}000 hours) that is \textbf{about 265 consumer GPUs}. This is a worst case: data-centre GPUs serving many pages at
once should need far fewer machines, which has to be measured. \emph{Training:} 1--2 data-centre GPUs with about 80\,GB
memory for LoRA fine-tuning of 8B--30B models. \emph{Storage:} about 40\,TB of scans a year (assuming about 0.5\,MB per
page), plus backups. \emph{Operations:} a reviewer web app, logins (examiner IDs already exist for e-marking), and a
check every season that models still work well.

\paragraph{7. What makes it different.} It checks the humans instead of replacing them, which matches regulators' current
position \citep{ofqual2026}; it checks every script, not a random sample; it is built for handwritten, bilingual scripts
that are often still on paper; simple, explainable rule flags come first and AI disagreement second; and it runs inside
the building on open models.

\paragraph{8. What I would build first.} I have already built the first version (Chapter~\ref{sec:prototype}). Next, a
\textbf{trial with one board, one subject and one exam season}: run the audit next to the existing rechecking, without
changing any marks, and compare what each finds.
